\section{Discussion}
\label{sec:discussion}

\subsection{Key Findings Interpretation}

\textbf{The large efficiency advantage as a practical benchmark.} GNN-NFN reaches 90\% of its peak R$^2$ at $N{\approx}147$ training models; flat-MLP requires the full dataset (${\sim}7{,}000$ models) to reach 90\% of its own peak R$^2{=}0.886$. An organization collecting model checkpoints for downstream weight-space analysis should target at least 250 models to benefit from equivariant encoders, and can achieve strong relative performance with as few as 150 models---while flat-MLP remains in early scaling territory.

\textbf{The crossover as a new constraint on equivariant inductive bias.} At $N{=}100$, GNN-NFN achieves below-chance prediction (R$^2{=}{-}0.016$)---worse than predicting the mean. Equivariant inductive bias was expected to help most at low data, where structural priors reduce the hypothesis search space. We hypothesize that the graph encoder architecture requires a minimum diversity of weight-graph topologies to learn useful node embeddings---below this threshold, the structural constraint provides no useful signal and may actively harm performance. However, LR miscalibration at small $N$ remains an unruled-out alternative explanation. An LR sweep at $N{=}100$ is the highest-priority ablation to distinguish these explanations. This connects to a broader pattern: specialist architectures require sufficient data to demonstrate their advantages. CNNs outperform MLPs on images when data is plentiful; at very small image datasets, the convolutional bias has insufficient examples to calibrate. Our results identify an analogous threshold for equivariant graph encoders on weight spaces.

\textbf{PermAug and structural equivariance as distinguishable strategies.} The $N{=}100$ crossover refutes the assumption that the two strategies are interchangeable. PermAug's benefit at $N{=}100$ appears primarily due to data quantity ($11{\times}$ expansion per gradient step); GNN-NFN's benefit at $N{\geq}250$ appears due to structural hypothesis-space constraint. The $N{=}500$--1000 convergence of the fraction toward 0.5--0.8 suggests both mechanisms are partially substitutable at intermediate scales.

\textbf{Connection to expressivity theory.} Full-scale convergence ($\Delta{=}0.008$ R$^2$) provides empirical support for \citet{Dayan2026Expressive}, and characterizes the \emph{window} of equivariant advantage: $N{\in}[150, 7000]$. The theoretical result establishes the ceiling parity; our data establish the practical efficiency window.

\subsection{Limitations}

\textbf{CIFAR-10 only.} All property-prediction efficiency results are from the CIFAR-10 CNN zoo. The MNIST MLP zoo---appropriate for DWSNets---was not available locally. The core efficiency claim is confirmed on CIFAR-10 only.

\textbf{DWSNets excluded from property prediction.} DWSNets requires $M{>}2$ FC layers; the CIFAR-10 CNN zoo has 2, making it architecturally incompatible. Equivariance is confirmed structurally ($\text{max\_diff}{=}7.45{\times}10^{-9}$) but property-prediction results for DWSNets are not reported.

\textbf{N=100 crossover is single-seed.} The most novel finding lacks multi-seed confidence intervals. We present it as a preliminary observation requiring 10-seed replication before claiming robustness.

\textbf{PermAug from separate experimental phase.} As described in Section~\ref{sec:method}, PermAug results come from a follow-up experiment sharing baselines from the primary experiment but using different random seeds. The crossover finding should be confirmed under a fully co-trained single-experiment design.

\textbf{Low zoo diversity.} CIFAR-10 zoo accuracy variance$=0.025$ (models cluster near convergence). Absolute R$^2$ values may be inflated relative to harder, more diverse zoos; the relative efficiency advantage is preserved.

\textbf{Hyperparameter fairness at small N.} Adam hyperparameters were tuned on full data. At $N{=}100$, the learning rate may disadvantage GNN-NFN's more complex architecture. An LR sweep at $N{=}100$ is needed to isolate this confound and distinguish structural limitation from training instability.

\textbf{No formal CI on efficiency ratio.} The interpolated $N_{\text{equiv},90}{\approx}147$ spans between $N{=}100$ and $N{=}250$, meaning the true $N_{\text{equiv},90}$ could range from 101 to 249. Using $N_{\text{plain},90}{=}7{,}000$ (full dataset), this yields efficiency ratios from ${\sim}28{\times}$ to ${\sim}69{\times}$. The central estimate of 47$\times$ is a point estimate under linear interpolation.

\subsection{Broader Impact}

This work provides actionable guidance for practitioners designing weight-space learning pipelines under data constraints. The decision boundary (${\sim}150$--250 models) reduces guesswork in encoder selection and zoo design. Reducing the data requirement for weight-space property prediction may lower barriers to applying these techniques in resource-constrained settings. There are no foreseeable negative societal impacts from this methodological work.
